\chapter{Introdução}
\label{cap:intro}

As Imagens Hiperespectrais (HSIs) constituem uma tecnologia de sensoriamento remoto que captura informações em centenas de bandas espectrais estreitas e contíguas, permitindo a identificação precisa de materiais e compostos presentes em uma cena \cite{Li2019DeepLF, Ahmad_2025}. Sua classificação apresenta desafios relacionados à alta dimensionalidade dos dados e à redundância espectral entre bandas adjacentes \cite{Qian2020, Lurs12162659}. As CNNs unidimensionais (CNNs 1D), que convoluem apenas o espectro de cada pixel, consomem menos memória e menos tempo por amostra que as bidimensionais e tridimensionais, ao custo de alguma acurácia, e são uma alternativa para cenários com restrições computacionais \cite{Viel10064297}.

A fusão de dados espaciais e espectrais melhora a classificação de HSIs \cite{IMANI202059, Souza2023}. \citeonline{Souza2023} integrou a uma CNN 1D o identificador do superpixel produzido pelo algoritmo SEEDS, concatenado às características extraídas pelas convoluções, e reportou, em Indian Pines, ganhos de 8,33, 8,79 e 9,47 pontos percentuais em OA, AA e Kappa em relação à mesma rede sem a informação espacial. O trabalho, porém, foi entregue como um protótipo: um notebook executado de forma sequencial, com hiperparâmetros ajustados manualmente e um único algoritmo de segmentação testado.

O processamento paralelo é uma estratégia consolidada para acelerar aplicações de aprendizado profundo \cite{Aach2025, Nandan10910577}, e as GPUs, com milhares de núcleos de processamento, executam simultaneamente operações independentes, como os elementos de um mapa de características de uma convolução \cite{CUDA_Programming, Jorda_cuDNN}. Na arquitetura de \citeonline{Souza2023} há ainda uma independência de outra natureza: a segmentação espacial e a extração de características espectrais só se encontram no ponto de fusão e poderiam ser executadas concorrentemente.

Este trabalho parte, portanto, da arquitetura de \citeonline{Souza2023} e a otimiza em quatro eixos: a organização do código, o paralelismo, a automação do ajuste de hiperparâmetros e a avaliação de outros algoritmos de segmentação. Os três últimos correspondem a trabalhos futuros sugeridos pelo próprio autor \cite{Souza2023}; o primeiro é pré-requisito para que os ganhos dos demais sejam atribuíveis. O ganho de cada eixo é medido contra uma referência própria, sempre com a fusão SEEDS.

\section{Problematização}
\label{sec:problematizacao}

A arquitetura de \citeonline{Souza2023} apresenta três limitações. A primeira é o processamento sequencial: as etapas de segmentação e classificação são executadas em série mesmo onde não dependem uma da outra. A segunda é o ajuste manual dos hiperparâmetros da rede e do segmentador: o próprio autor relata ter alterado, por conjunto de dados, a paciência da parada antecipada e outros parâmetros de treino \cite{Souza2023}, e cada nova cena exige repetir esse ajuste à mão. A terceira é que o algoritmo de segmentação foi escolhido sem comparação com alternativas.

Há um problema anterior a esses três: para que qualquer otimização seja atribuível, é preciso reproduzir o trabalho original de forma auditável. Um notebook com estado global e sem semente fixada para o treino da rede não permite afirmar que uma diferença medida vem da otimização, e não de uma execução diferente (\autoref{sec:resLinhaBase}).

\subsection{Solução Proposta}
\label{subsec:solucao}

A solução proposta consiste em reimplementar o trabalho de \citeonline{Souza2023} em um pacote Python modular e determinístico, que preserve a arquitetura da rede, o PCA, a divisão dos dados e os hiperparâmetros de treino do original, e, sobre essa base, implementar e medir: (\textit{i}) estratégias de paralelismo entre os ramos do pré-processamento, por partição da cena e na fusão concorrente; (\textit{ii}) a busca automática de hiperparâmetros da CNN e do segmentador; e (\textit{iii}) a comparação de quatro algoritmos de segmentação sob otimização conjunta.

\section{Objetivos}

\subsection{Objetivo Geral}

Analisar o impacto de técnicas de segmentação, parametrização automática e processamento paralelo na classificação de imagens hiperespectrais utilizando redes neurais convolucionais unidimensionais.

\bigskip
\subsection{Objetivos Específicos}

\begin{enumerate}
    \item Reimplementar a arquitetura de \citeonline{Souza2023} de forma modular e reprodutível, preservando o PCA, a divisão dos dados e os hiperparâmetros de treino da implementação original e garantindo resultados idênticos entre execuções com a mesma configuração;

    \item Explorar algoritmos de segmentação espacial (SEEDS, SLIC, LSC, Watershed) para o processamento de HSIs e avaliar seu impacto na classificação;

    \item Implementar métodos de ajuste automático de hiperparâmetros (busca aleatória, Otimização Bayesiana e algoritmo genético) para a arquitetura de classificação;

    \item Implementar estratégias de processamento paralelo para acelerar as etapas de segmentação e classificação espectral, explorando a independência de dados entre elas e a execução concorrente via threads e processos;

    \item Disponibilizar uma análise comparativa do fluxo de classificação de HSIs baseando-se em métricas de acurácia (OA, AA, Kappa) e custo computacional.
\end{enumerate}

\section{Metodologia}

Quanto aos objetivos, a pesquisa é exploratória e descritiva \cite{Gil2002}. A abordagem é quantitativa, e as conclusões foram obtidas de maneira indutiva, a partir da comparação das métricas de desempenho entre configurações.

Foi realizada uma pesquisa bibliográfica sobre aprendizado profundo \cite{DeeplearningLC, DeepLearningIG}, processamento de imagens hiperespectrais \cite{Li2019DeepLF, Ahmad_2025, Qian2020}, redes neurais convolucionais \cite{Gradient-basedlearningappliedtodocumentrecognition, NIPS2012_c399862d}, otimização de hiperparâmetros \cite{WU201926, liashchynskyi2019gridsearchrandomsearch, Bergstra2012Random, Bergstra2011TPE}, segmentação de imagens \cite{bergh2013seedssuperpixelsextractedenergydriven, Achanta2012SLIC, Zhengqin7298741, Najman546254}, concorrência \cite{Silberschatz2018, Amdahl1967} e paralelização computacional \cite{Aach2025, Nandan10910577, CUDA_Programming}. Em seguida, o trabalho original foi reimplementado em PyTorch, e as técnicas escolhidas foram implementadas sobre essa base. A avaliação foi feita por uma campanha experimental com três conjuntos de dados de referência (Indian Pines, University of Pavia e Salinas Valley), organizada em experimentos retomáveis cujos resultados ficam registrados em arquivos versionados.

\section{Estrutura do Trabalho}

O \autoref{cap:fund} apresenta a fundamentação teórica e os trabalhos relacionados. O \autoref{cap:proj} define o projeto, com os requisitos e o plano de verificação. O \autoref{cap:desenvolvimento} descreve a implementação e os desvios em relação ao projeto. O \autoref{cap:resultados} expõe os resultados, organizados pelos quatro eixos de otimização, e a verificação dos requisitos. Por fim, o \autoref{cap:consid} retoma os objetivos e as hipóteses e aponta trabalhos futuros.
